% 和文 1 段組で，句読点を **`．，`（全角ピリオド・全角コンマ）** で組んだ論文．
%
% 情報処理学会・電子情報通信学会の標準表記で，和文論文の相当な割合がこれ．
% `．，` しか知らない文分割では**節まるごとが 1 文**になる（実測: 3,648 字・矩形 96 個の
% カードが 1 枚，画面には「0 / 1 文」）．さらに長すぎる文の上限が発火すると，
% 和文には語間の空白が無いので `…到達可能な処` / `理速度の…` と語の途中で切れる．
%
% japanese-onecolumn.tex と本文は同じで，句読点だけを入れ替えてある．
% 2 本を並べて測れば，「表記だけが違う同じ原稿」で結果が変わらないことを確かめられる．
% ----------------------------------------------------------------------------
% 計測（npm run measure）への足し方
%
%   1. この .tex と .pdf と .truth.json を scripts/real-papers/ に置く．
%      measureRealPapers() は real-papers/*.pdf を総当たりで読むので，
%      measure-quality.ts 側に足す記述は要らない．
%   2. 組み直しには **xelatex** が要る（和文の書体を使うため pdflatex では通らない）．
%      scripts/build-real-papers.ts は pdflatex 決め打ちなので，
%      「\usepackage{fontspec} があれば xelatex」のような分岐を足す．
%      正解（.truth.json）の作り方は欧文と同じ scripts/real-paper-truth.ts のまま
%      で通る（和文の \s{...} もそのまま拾える）．
%   3. verifyInlineMath（pdftotext との突き合わせ）は \m{...}{...} を 1 つも
%      使っていないので，この原稿では何もしない．
% ----------------------------------------------------------------------------
\documentclass[11pt,a4paper]{article}

\usepackage{fontspec}
\setmainfont{Noto Serif CJK JP}
\setmonofont{DejaVu Sans Mono}
\XeTeXlinebreaklocale "ja"
% 和文の字間に入れる伸縮．pTeX の \kanjiskip（10pt で plus/minus 0.4pt）に合わせる．
% ここを大きく取ると，pdf.js が字間の空きを空白文字と見なして
% （TRACKING_SPACE_FACTOR = 0.102 em）文字片の中に空白を撒く．実在の和文組版は
% そこまで空けないので，現実的な版面にしてから測る．
\XeTeXlinebreakskip=0pt plus 0.4pt minus 0.4pt
\usepackage[margin=2.4cm]{geometry}
\usepackage{fancyhdr}
\usepackage{amsmath}

\newcommand{\s}[1]{#1}
\renewcommand{\abstractname}{概要}
% 参考文献の見出しは既定の "References" のままにする．
% layout.ts の KNOWN_HEADINGS / TAIL_HEADINGS は欧文の見出し語しか知らないので，
% ここを「参考文献」にすると参考文献の節が丸ごと本文として扱われる．
% それは文の連結とは別の問題なので，ここでは切り分けて測る．

\pagestyle{fancy}
\fancyhf{}
\fancyhead[L]{和文プレプリント}
\fancyfoot[C]{\thepage}

\title{日本語の本文を含む論文における文境界の推定}
\author{坂比賀 悠真\\ 山手大学 情報科学研究科}
\date{}

\begin{document}
\maketitle

\begin{abstract}
\s{本稿では，日本語の本文に英語の術語が混ざる論文を対象に，段落と文の切り出しを検討する．}%
\s{既存の手法は空白で語を区切ることを前提としており，日本語の本文では文の境界を見失う．}%
\s{我々は句点と括弧の対応を手掛かりとする単純な規則を提案し，三つのコーパスで評価した．}%
\end{abstract}

\section{はじめに}

\s{学術論文の多くは英語で書かれているが，国内の学会誌には日本語で書かれたものも多い．}%
\s{そうした論文では，術語だけが英語のまま残されることが普通である．}%
\s{たとえば attention や transformer といった語は，訳さずにそのまま使われる．}%
\s{この混在は，空白を手掛かりに語を切る解析器にとって厄介な性質を持つ．}%
\s{英語の部分では空白が語の区切りになるが，日本語の部分では空白がまったく現れない．}%
\s{行の折り返しは語の途中でも起きるため，行を素朴に繋ぐと語が割れてしまう．}%

\s{本稿の貢献は三つある．}%
\s{第一に，日本語と英語が混ざった本文で文の境界を求める規則を示す．}%
\s{第二に，その規則を三つのコーパスで評価し，既存手法との差を報告する．}%
\s{第三に，残った誤りを分類し，どこに限界があるかを明らかにする．}%

\section{関連研究}

\s{文の境界推定は古くから研究されており，英語では終止符と大文字の組み合わせが主な手掛かりである．}%
\s{日本語では句点がほぼ唯一の手掛かりだが，括弧の中の句点を数えてしまう誤りが多い．}%
\s{実装としては \texttt{nltk} や \texttt{spacy} が広く使われているが，いずれも空白で語が区切られることを前提にしている．}%
\s{近年は学習に基づく手法も提案されているが，分野ごとの学習データが必要になる \cite{sato2021}．}%
\s{本稿が対象とするのは，学習データを用意できない場面である．}%

\section{提案手法}

\s{提案手法は，句点の直後で切ることを原則とし，括弧の対応が取れていない位置では切らない．}%
\s{英語の術語が括弧の中に置かれる場合を考慮し，全角と半角の括弧を同じものとして扱う．}%
\s{また，箇条書きの行頭記号は文の始まりとみなし，直前で必ず切る．}%
\s{入力の長さを $n$ とすると，提案手法の計算量は文字数に比例し，追加の記憶領域は定数である．}%
\s{規則はすべて決定的であり，同じ入力に対して必ず同じ出力を返す．}%

\begin{figure}[t]
  \centering
  \rule{6cm}{2.4cm}
  \caption{提案手法の処理の流れ．行の継ぎ目で空白を入れるかどうかを文字の種類で決める．}
\end{figure}

\section{実装}

\s{提案手法は五百行程度の手続き型の実装であり，外部の辞書も学習済みの模型も使わない．}%
\s{入力は版面から取り出した文字片の列で，出力は文の列である．}%
\s{文字片は座標を持つので，行にまとめる処理を先に通してから文に割る．}%
\s{行にまとめる処理は，段の切れ目を空白の帯として検出し，左の段から順に読む．}%
\s{この順序を誤ると，左右の段が一行の中で混ざり，文がすべて壊れる．}%

\s{行を繋ぐ位置では，継ぎ目の両側の文字を見て空白の要否を決める．}%
\s{両側が和文であれば空白を入れず，両側が欧文であれば空白を入れる．}%
\s{片側だけが和文である場合は，同じ文書の行の途中の書き方に合わせる．}%
\s{この判断は文書ごとに一度だけ行い，以降は同じ判断を使い回す．}%
\s{判断の材料が一つも無い文書では，欧文と同じ扱いにする．}%

\section{評価}

\s{三つのコーパスは，情報処理学会論文誌，人工知能学会誌，および arXiv の和文要旨から構成される．}%
\s{いずれも人手で文の境界を付与し，二人の作業者の一致率は九十七パーセントであった．}%
\s{提案手法の $F$ 値は三つのコーパスでそれぞれ九十四，九十二，八十九であった．}%
\s{既存手法との差はいずれのコーパスでも有意であった．}%
\s{誤りの多くは，箇条書きと表の見出しが本文に混ざる場合に集中していた．}%

\s{処理時間は百ページの論文で一秒未満であり，実用上の問題は無い．}%
\s{記憶領域は文字片の数に比例し，百ページの論文で数十メガバイトに収まる．}%
\s{いずれの指標も既存の実装と同等であり，速度を犠牲にしていない．}%

\section{考察}

\s{残った誤りの半分は，見出しが本文として扱われたことに起因する．}%
\s{和文の見出しは番号と本文の間に空白が一つ入るだけで，欧文のような大文字の手掛かりが無い．}%
\s{文字の大きさだけを手掛かりにすると，本文と同じ大きさで組まれた小見出しを取り逃がす．}%
\s{この問題は文の切り出しとは独立であり，行の役割を決める処理の側で解くべきである．}%

\s{もう一つの誤りは，柱と本文が同じ段に並ぶ場合に起きる．}%
\s{柱は複数のページに同じ字面で現れるので，ページ数が少ない文書では判別が難しい．}%
\s{この場合も文の切り出しではなく，行の役割を決める処理の問題である．}%

\section{誤りの分析}

\s{誤りを種類ごとに数えると，見出しに起因するものが最も多く，次に柱に起因するものが続いた．}%
\s{見出しに起因する誤りは，見出しと直後の本文が一枚の札に同居する形で現れる．}%
\s{柱に起因する誤りは，ページの変わり目で本文の途中に柱の文字列が挟まる形で現れる．}%
\s{どちらも読む側からは「原文がそう書いてある」ようにしか見えないため，気付きにくい．}%
\s{気付きにくい誤りは，気付ける誤りよりも害が大きいと我々は考える．}%

\s{三つ目の種類は，図表の見出しが本文として扱われる誤りである．}%
\s{和文の図表の見出しは「図 1」「表 2」の形をとり，欧文の版面とは接頭辞が異なる．}%
\s{接頭辞の一覧に和文の語を加えれば解けるが，雑誌ごとに表記が揺れるため万全ではない．}%
\s{我々は，接頭辞に頼らず，版面の位置と前後の行との間隔から判断する方法を検討している．}%

\s{四つ目の種類は，箇条書きの各項目が一つの文として扱われる誤りである．}%
\s{箇条書きの項目は句点で終わらないことが多く，次の項目と繋がってしまう．}%
\s{行頭の記号を手掛かりにすれば切れるが，記号の種類は版面によって異なる．}%
\s{この種類の誤りは全体の一割に満たないため，優先度は低いと判断した．}%

\section{他の言語への拡張}

\s{中国語は日本語と同じく分かち書きをしないので，同じ規則がそのまま使える．}%
\s{一方，朝鮮語は分かち書きをするため，行の継ぎ目には欧文と同じく空白が要る．}%
\s{同じ東アジアの文字であっても，空白の扱いは言語によって異なる．}%
\s{したがって，文字の範囲を「東アジアの文字」とひとまとめにする実装は誤りである．}%
\s{我々の実装は，漢字・ひらがな・カタカナ・全角記号のみを対象とし，ハングルを含めない．}%

\s{縦書きの版面については，行の向きが変わるだけで規則そのものは変わらない．}%
\s{ただし，版面から行を組み立てる処理は縦書きを想定していないため，別に手当てが要る．}%
\s{振り仮名は本文と重なる位置に置かれるため，行の組み立ての段階で落とす必要がある．}%
\s{これらは今後の課題である．}%

\section{おわりに}

\s{日本語と英語が混ざった論文の文境界推定について，学習を必要としない規則を示した．}%
\s{三つのコーパスで既存手法を上回る結果が得られた．}%
\s{今後は，数式が本文に混ざる場合への拡張を検討したい．}%

\begin{thebibliography}{9}
% 参考文献は欧文で書く．和文の論文でも引用先は英語文献であることが多く，
% そのほうが実物に近い．加えて layout.ts の参考文献の判定
% （REFERENCE_ENTRY = /^\[\d+\]\s*\p{Lu}/u）は「[1] の後ろが**欧文の大文字**」を
% 求めるので，"[1] 佐藤 花子" 形式の和文の参考文献はいまは本文に混ざる．
% それは文の連結とは別の問題なので，ここでは切り分けて測る．
\bibitem{sato2021}
H. Sato and I. Suzuki. Sentence boundary detection for Japanese documents revisited. Journal of Information Processing, 62(3):512--523, 2021.
\bibitem{tanaka2019}
A. Tanaka. Structural analysis of scholarly documents. Journal of the Japanese Society for Artificial Intelligence, 34(2):88--97, 2019.
\end{thebibliography}

\end{document}
